% Generated draft: requires mathematical and source review.
\documentclass[11pt]{article}
\usepackage[margin=25mm]{geometry}
\usepackage{fontspec}
\setmainfont{Latin Modern Roman}
\usepackage{amsmath,amssymb,mathtools,unicode-math}
\setmathfont{Latin Modern Math}
\usepackage{enumitem,needspace,xurl,hyperref}
\setlength{\parindent}{0pt}
\setlength{\parskip}{5pt}
\setlength{\emergencystretch}{3em}
\newcommand{\N}{\mathbb{N}}
\newcommand{\Z}{\mathbb{Z}}
\newcommand{\Q}{\mathbb{Q}}
\newcommand{\R}{\mathbb{R}}
\newcommand{\C}{\mathbb{C}}
\newcommand{\F}{\mathbb{F}}
\newcommand{\A}{\mathbb{A}}
\newcommand{\PP}{\mathbb{P}}
\newcommand{\E}{\mathbb{E}}
\newcommand{\eps}{\varepsilon}
\newcommand{\dd}{\,\mathrm{d}}
\begin{document}
\section*{UnsolvedMath ID: 20002900}
\textbf{Source problem code:} catalogue entry
\subsection{Definitions and mathematical statement}
Fix \(k\ge1\), put \(n=k+1\), and let \(\Lambda\) be the integral graded ring of symmetric functions, with complete functions \(h_i\), monomial functions \(m_\lambda\), and the Hall pairing. Set
\[
\Lambda_{(k)}=\Z[h_1,\ldots,h_k],\qquad
\Lambda^{(k)}=\Lambda/\langle m_\lambda:\lambda_1>k\rangle.
\]
These are graded dual under the Hall pairing. Index their standard \(t=1\) \(k\)-Schur and dual \(k\)-Schur bases by \(k\)-bounded partitions \(\lambda\), meaning \(\lambda_1\le k\): write them \(s_\lambda^{(k)}\) and \(\mathfrak S_\lambda^{(k)}\), with
\[
\langle s_\lambda^{(k)},\mathfrak S_\mu^{(k)}\rangle=\delta_{\lambda\mu}.
\]
Let \(\operatorname{Gr}=SL_n(\C((t)))/SL_n(\C[[t]])\) be the affine Grassmannian. Its Schubert classes are indexed by affine Grassmannian elements of the affine symmetric group, identified with \(k\)-bounded partitions by the standard bijection. Denote the homology and cohomology Schubert classes by \(\xi_\lambda\) and \(\xi^\lambda\), in topological degrees \(2|\lambda|\).

The source requests the type-\(A\) identifications
\[
\Lambda_{(k)}\xrightarrow{\sim}H_*(\operatorname{Gr};\Z),
\quad s_\lambda^{(k)}\longmapsto\xi_\lambda,
\qquad
\Lambda^{(k)}\xrightarrow{\sim}H^*_{\rm gr}(\operatorname{Gr};\Z),
\quad\mathfrak S_\lambda^{(k)}\longmapsto\xi^\lambda.
\]
The homology multiplication is the Pontryagin product and cohomology uses cup product; \(H^*_{\rm gr}\) denotes the direct sum of the graded cohomology groups, dual degree by degree to homology. If a completed cohomology convention is used instead, the symmetric-function side must be completed as well. No extra deformation parameter \(t\) is retained in the \(k\)-Schur functions here.

The phrase ``symmetric component of affine Schubert polynomials'' is interpreted by the specialization in the affine-flag representatives mentioned in the supplied aid: an affine Schubert representative indexed by \(w\) specializes to its affine Stanley symmetric function \(F_w\), and for the affine Grassmannian index \(w_\lambda\),
\[
F_{w_\lambda}=\mathfrak S_\lambda^{(k)}
\]
as a class in the dual symmetric-function model. A precise polynomial representative ring and specialization map require the conventions of the cited work, not a literal unspecified decomposition into a symmetric and nonsymmetric summand.

The supplied aid attributes the homology and cohomology statements to Lam and the affine-flag representative interpretation to Lee, and reports them as solved in the literature. It also reports the following low-degree consequence: in combinatorial degrees \(d\le k\), both models have ordinary Schur representatives \(s_\lambda\) (in the quotient where appropriate), and products of total degree at most \(k\) have the usual Littlewood--Richardson coefficients. At degree \(k+1\), each affine model has rank \(p(k+1)-1\), where \(p(d)\) counts partitions of \(d\), whereas \(\Lambda_{k+1}\) has rank \(p(k+1)\). Hence the full uniform identification with unrestricted ordinary Schur calculus cannot extend to that degree; this does not preclude individual higher-degree coincidences.
\subsection{Short English statement}
Identify the type-\(A\), \(t=1\) \(k\)-Schur and dual bases with affine-Grassmannian homology and cohomology Schubert classes, and interpret the older symmetric-component phrase through affine Stanley specialization. The supplied report labels these literature results solved and records a sharp uniform low-degree comparison with ordinary Schur calculus.
\subsection{Sources}
\begin{itemize}
\item AIM workshop \emph{Generalized Kostka polynomials}, Problem (5), \url{https://aimath.org/WWN/kostka/prob2.pdf}.
\item The supplied research summary attributes the type-\(A\) affine-Grassmannian results to Lam and the affine-flag representative specialization to Lee. No full titles, publication dates, theorem numbers or bibliography for those works are supplied.
\item Local UnsolvedMath registry, permanent ID 20002900, exact code \texttt{catalogue entry}; canonical location \texttt{aim-special-functions-notes.json}, notes[4], dated 14 August 2026; statement labeled exact and registry status solved, research classification solved in literature.
\item Research attempt 1: partial result, problem solved at the run, rigor label proved, novelty confidence low. The research summary and low-degree consequence are machine-generated aids, not peer reviewed, requiring independent verification before citation.
\end{itemize}
\end{document}
